% ============================================================================
% FORMULÁRIO TÉCNICO, GLOSSÁRIO ALFABÉTICO E REFERÊNCIAS OFICIAIS (CHF662)
% ============================================================================
\clearpage
\phantomsection
\addcontentsline{toc}{chapter}{Formulário Técnico, Glossário e Referências Oficiais}
\chapter*{Formulário Técnico, Glossário e Referências Oficiais}
\markboth{FORMULÁRIO, GLOSSÁRIO E REFERÊNCIAS}{FORMULÁRIO, GLOSSÁRIO E REFERÊNCIAS}

Este capítulo de encerramento reúne em formato de consulta rápida as equações, escalas classificatórias, tabelas operacionais, verbetes conceituais e a bibliografia integral do componente curricular \textbf{CHF662, Análise Ambiental} (DCHF/UEFS). Utilize-o como guia de revisão para avaliações, atividades de campo no Semiárido e elaboração do Diagnóstico Ambiental Aplicado.

% ============================================================================
\section*{1. Formulário Técnico e Tabelas de Consulta Rápida}
\addcontentsline{toc}{section}{1. Formulário Técnico e Tabelas de Consulta Rápida}
% ============================================================================

\subsection*{(a) Risco Ambiental, AIA e Matriz Ponderada de Impactos (Caps. 4 e 5)}

\begin{center}
\renewcommand{\arraystretch}{1.28}
\small
\begin{tabularx}{\textwidth}{@{}p{3.4cm} p{4.6cm} >{\raggedright\arraybackslash}X@{}}
  \toprule
  \rowcolor{uefsNavy}
  \textcolor{white}{\bfseries Instrumento / Índice} & \textcolor{white}{\bfseries Expressão / Estrutura} & \textcolor{white}{\bfseries Significado e Critério de Interpretação} \\
  \midrule
  Cadeia Causal (Sánchez, 2013) & $\text{Atividade} \to \text{Aspecto} \to \text{Impacto}$ & \textit{Atividade}: ação humana; \textit{Aspecto}: mecanismo de interação com o meio; \textit{Impacto}: alteração (adversa ou benéfica) na qualidade ambiental. \\
  \rowcolor{bgGray}
  Risco Ambiental (Matriz $4 \times 4$) & $R = P \times V$ \newline ($P, V \in \{1, 2, 3, 4\}$) & $P$: Perigo/ameaça física; $V$: Vulnerabilidade de pessoas/bens. Classes: \textbf{Baixo} ($1$--$3$), \textbf{Médio} ($4$--$6$), \textbf{Alto} ($8$--$9$), \textbf{Crítico} ($12$--$16$). \\
  Índice de Significância Ponderado (AIA) & $I_S = \displaystyle\sum_{j=1}^{k} w_j \cdot n_j$, \quad $\displaystyle\sum_{j=1}^{k} w_j = 1$ & $w_j$: peso do critério $j$ (magnitude, abrangência, duração, reversibilidade); $n_j$: nota padronizada (ex.: $1$ a $4$). Permite hierarquizar impactos de forma auditável. \\
  \rowcolor{bgGray}
  Hierarquia de Mitigação & $\text{Evitar} > \text{Minimizar} > \text{Recuperar} > \text{Compensar}$ & Ordem obrigatória de prioridade na AIA: prevenir na origem (alternativa locacional/tecnológica) antes de mitigar, reabilitar ou compensar impactos residuais. \\
  \bottomrule
\end{tabularx}
\end{center}

\subsection*{(b) Morfometria de Bacias Hidrográficas e Declividade de Vertentes (Caps. 7 e 8)}

\begin{center}
\renewcommand{\arraystretch}{1.30}
\small
\begin{tabularx}{\textwidth}{@{}p{3.3cm} p{4.5cm} >{\raggedright\arraybackslash}X@{}}
  \toprule
  \rowcolor{uefsNavy}
  \textcolor{white}{\bfseries Parâmetro} & \textcolor{white}{\bfseries Fórmula Operacional} & \textcolor{white}{\bfseries Unidades e Interpretação Hidrogeomorfológica} \\
  \midrule
  Densidade de Drenagem ($D_d$) & $D_d = \dfrac{L_t}{A}$ & $\text{km/km}^2$ ($L_t$: soma dos comprimentos dos canais; $A$: área). Baixa ($< 0,5$), média ($0,5$--$2,0$), alta ($> 2,0$): reflete impermeabilidade/dissecação. \\
  \rowcolor{bgGray}
  Coeficiente de Compacidade ($K_c$) & $K_c = 0,28 \cdot \dfrac{P}{\sqrt{A}}$ & Adimensional ($P$: perímetro em km; $A$ em $\text{km}^2$). Sempre $K_c \ge 1,0$. Próximo de $1,0$: bacia circular (maior risco de cheias rápidas); $K_c > 1,5$: alongada. \\
  Índice de Circularidade ($I_c$) & $I_c = 12,57 \cdot \dfrac{A}{P^2}$ & Adimensional ($0 < I_c \le 1,0$). Valores $> 0,51$ indicam tendência circular e picos de vazão concentrados; valores baixos indicam bacia alongada. \\
  \rowcolor{bgGray}
  Amplitude Altimétrica ($H_m$ ou $\Delta H$) & $H_m = H_{\max} - H_{\min}$ & $\text{m}$. Diferença entre a cota mais elevada do divisor e a cota do exutório; expressa o potencial energético gravitacional da bacia. \\
  Relação de Relevo ($R_r$) & $R_r = \dfrac{H_{\max} - H_{\min}}{L_a}$ & $\text{m/km}$ ou $\text{m/m}$ ($L_a$: comprimento axial da bacia). Quanto maior $R_r$, maior a declividade média global e a energia erosiva. \\
  \rowcolor{bgGray}
  Hierarquia de Strahler e Razão de Bifurcação ($R_b$) & $1+1\to 2$; $2+2\to 3$; $2+1\to 2$ \newline $R_b = \dfrac{N_\omega}{N_{\omega+1}}$ & $N_\omega$: número de segmentos de ordem $\omega$. Valores típicos naturais: $R_b \approx 3,0$ a $5,0$ (valores maiores indicam forte controle estrutural/relevo). \\
  Declividade de Vertente ($S\%$ e $\theta^\circ$) & $S(\%) = \dfrac{\Delta H}{L_h}\times 100$ \newline $\theta = \arctan\!\left(\dfrac{\Delta H}{L_h}\right)$ & $\Delta H$: desnível vertical (m); $L_h$: distância horizontal (m). Atenção: $S = 100\% \iff \theta = 45^\circ$ (nunca confundir $\%$ com graus!). \\
  \bottomrule
\end{tabularx}
\end{center}

\noindent\textbf{Classes de Relevo (EMBRAPA) por declividade $S(\%)$:} Plano ($0$--$3\%$); Suave ondulado ($3$--$8\%$); Ondulado ($8$--$20\%$); Forte ondulado ($20$--$45\%$); Montanhoso ($45$--$75\%$); Escarpado ($> 75\%$).

\subsection*{(c) Ecodinâmica de Tricart (1977) e Fragilidade Ambiental de Ross (1994) (Cap. 9)}

\begin{center}
\renewcommand{\arraystretch}{1.22}
\small
\begin{tabularx}{\textwidth}{@{}p{2.9cm} p{3.6cm} >{\raggedright\arraybackslash}X@{}}
  \toprule
  \rowcolor{uefsNavy}
  \textcolor{white}{\bfseries Meio (Tricart)} & \textcolor{white}{\bfseries Balanço Dominante} & \textcolor{white}{\bfseries Características e Implicações para o Uso} \\
  \midrule
  \textbf{Estável} & $\text{Pedogênese} > \text{Morfogênese}$ & Cobertura vegetal fechada, relevo suave, solos evoluídos e baixa taxa de remoção superficial; admite usos conservacionistas planejados. \\
  \rowcolor{bgGray}
  \textbf{Intergrade} & $\text{Pedogênese} \approx \text{Morfogênese}$ & Transição dinâmica sensível a perturbações antrópicas; manejo inadequado desloca rapidamente o sistema para instabilidade. \\
  \textbf{Fortemente Instável} & $\text{Morfogênese} > \text{Pedogênese}$ & Relevo íngreme, solos rasos ou solo exposto, chuvas torrenciais e erosão/movimentos de massa ativos; exige preservação ou recuperação prioritária. \\
  \bottomrule
\end{tabularx}
\end{center}

\vspace{4pt}
\begin{center}
\renewcommand{\arraystretch}{1.20}
\small
\begin{tabularx}{\textwidth}{@{}c p{2.3cm} p{2.3cm} p{3.5cm} >{\raggedright\arraybackslash}X@{}}
  \toprule
  \rowcolor{uefsTeal}
  \textcolor{white}{\bfseries Classe} & \textcolor{white}{\bfseries Fragilidade (Ross)} & \textcolor{white}{\bfseries Declividade} & \textcolor{white}{\bfseries Solos Típicos (Exemplos)} & \textcolor{white}{\bfseries Grau de Proteção do Uso da Terra} \\
  \midrule
  \textbf{1} & Muito Fraca & $< 6\%$ & Latossolos muito profundos & Mata densa / floresta preservada (proteção muito alta) \\
  \rowcolor{bgGray}
  \textbf{2} & Fraca & $6\%$ a $12\%$ & Latossolos / Argissolos prof. & Capoeira / mata secundária / silvicultura densa \\
  \textbf{3} & Média & $12\%$ a $20\%$ & Argissolos / Planossolos & Pastagens bem manejadas / culturas perenes \\
  \rowcolor{bgGray}
  \textbf{4} & Forte & $20\%$ a $30\%$ & Cambissolos / solos médios & Cultivos anuais / pastagens degradadas \\
  \textbf{5} & Muito Forte & $> 30\%$ & Neossolos Litólicos / Quartzarênicos & Solo exposto / terraplenagem / mineração (proteção muito baixa) \\
  \bottomrule
\end{tabularx}
\end{center}

\noindent\textbf{Matriz de Cruzamento de Ross ($FP \times \text{Uso} \to FE$):} $FP$ sintetiza os fatores físicos (declividade, solos, clima); o cruzamento com a classe de proteção do uso da terra gera a Fragilidade Emergente ($FE$). Exemplo: $FP = 4$ sob mata densa (classe 1) atenua-se para $FE = 2$ (Fraca), mas sob solo exposto (classe 5) agrava-se para $FE = 5$ (Muito Forte).

\subsection*{(d) Indicadores de Qualidade do Solo e Equação Universal de Perda de Solo: USLE (Caps. 11 e 12)}

\begin{center}
\renewcommand{\arraystretch}{1.30}
\small
\begin{tabularx}{\textwidth}{@{}p{3.4cm} p{4.8cm} >{\raggedright\arraybackslash}X@{}}
  \toprule
  \rowcolor{uefsNavy}
  \textcolor{white}{\bfseries Indicador / Modelo} & \textcolor{white}{\bfseries Equação} & \textcolor{white}{\bfseries Parâmetros, Unidades e Referências} \\
  \midrule
  Densidade do Solo ($D_s$) & $D_s = \dfrac{m_s}{V}$ & $m_s$: massa de solo seco a $105\,^\circ\text{C}$ (g); $V$: volume do anel ($\text{cm}^3$). Em solos minerais, $D_s > 1,55\text{--}1,65\ \text{g/cm}^3$ indica compactação crítica. \\
  \rowcolor{bgGray}
  Porosidade Total ($P_t$) & $P_t(\%) = \left(1 - \dfrac{D_s}{D_p}\right)\times 100$ & $D_p$: densidade de partículas (padrão mineral $\approx 2,65\ \text{g/cm}^3$). Solos estruturados apresentam $P_t \approx 45\%$ a $60\%$. \\
  Conversão $CO \to MO$ & $MO = 1,724 \times CO$ & Fator de Van Bemmelen (assume $\approx 58\%$ de carbono na matéria orgânica). Expresso em $\text{g/kg}$ ou $\%$ ($1\% = 10\ \text{g/kg}$). \\
  \rowcolor{bgGray}
  Normalização Linear (IQS: $0$ a $1$) & \textbf{Maior melhor:} $S_i = \dfrac{x - x_{\min}}{x_{\max} - x_{\min}}$ \newline \textbf{Menor melhor:} $S_i = \dfrac{x_{\max} - x}{x_{\max} - x_{\min}}$ & Usada para compor o Índice de Qualidade do Solo ($\text{IQS} = \sum w_i S_i$ ou média simples). \textit{Maior melhor}: $MO, P_t$, infiltração; \textit{Menor melhor}: $D_s$, resistência à penetração. \\
  Equação Universal de Perda de Solo (USLE) & $A = R \cdot K \cdot L \cdot S \cdot C \cdot P$ & $A$: perda de solo ($\text{t}\cdot\text{ha}^{-1}\cdot\text{ano}^{-1}$); $R$: erosividade ($\text{MJ}\cdot\text{mm}\cdot\text{ha}^{-1}\cdot\text{h}^{-1}\cdot\text{ano}^{-1}$); $K$: erodibilidade ($\text{t}\cdot\text{h}\cdot\text{MJ}^{-1}\cdot\text{mm}^{-1}$); $L, S, C, P$: adimensionais (comprimento, declividade, cobertura/manejo e práticas conservacionistas). Comparar $A$ com a tolerância $T$ (geralmente $5$ a $12\ \text{t}\cdot\text{ha}^{-1}\cdot\text{ano}^{-1}$). \\
  \bottomrule
\end{tabularx}
\end{center}

\subsection*{(e) Bioengenharia de Água e Solo: Estabilidade de Taludes e Paliçadas (Cap. 13)}

\begin{center}
\renewcommand{\arraystretch}{1.30}
\small
\begin{tabularx}{\textwidth}{@{}p{3.4cm} p{4.8cm} >{\raggedright\arraybackslash}X@{}}
  \toprule
  \rowcolor{uefsNavy}
  \textcolor{white}{\bfseries Equação / Critério} & \textcolor{white}{\bfseries Expressão Matemática} & \textcolor{white}{\bfseries Variáveis e Aplicação Prática} \\
  \midrule
  Mohr--Coulomb com Coesão Radicular & $\tau_f = (c' + c_r) + (\sigma - u)\tan\phi'$ & $\tau_f$: resistência ao cisalhamento (kPa); $c'$: coesão efetiva do solo; $c_r$: acréscimo de coesão por raízes (kPa); $\sigma$: tensão normal; $u$: poropressão; $\phi'$: ângulo de atrito interno. \\
  \rowcolor{bgGray}
  Fator de Segurança ($FS$) & $FS = \dfrac{\text{Forças resistentes}}{\text{Forças solicitantes}}$ & $FS < 1,0$: ruptura iminente; $FS = 1,0$: equilíbrio limite; projeto estável exige $FS \ge 1,3$ a $1,5$ (elevado por $c_r$ e redução de $u$). \\
  Espaçamento de Paliçadas Permeáveis & $E = \dfrac{h}{S}$ \qquad e \qquad $n = \dfrac{L}{E}$ & $h$: altura útil do vertedouro (m); $S$: declividade longitudinal do canal em $\text{m/m}$ (ex.: $10\% = 0,10$); $L$: comprimento do trecho da ravina (m). \\
  \rowcolor{bgGray}
  Volume de Sedimento Retido (Cunha) & $V_{\text{cunha}} = \dfrac{1}{2}\,b\,h\,E$ \quad ($\text{m}^3$) & $b$: largura média da seção do canal (m); $h$: altura útil (m); $E$: comprimento da cunha retida (m). Capacidade total em série: $V_{\text{total}} = n \cdot V_{\text{cunha}}$. \\
  \bottomrule
\end{tabularx}
\end{center}

\subsection*{(f) Resumo Operacional do Código Florestal (Lei nº 12.651/2012) e Licenciamento (Cap. 14)}

\begin{center}
\renewcommand{\arraystretch}{1.22}
\small
\begin{tabularx}{\textwidth}{@{}p{4.3cm} p{4.0cm} >{\raggedright\arraybackslash}X@{}}
  \toprule
  \rowcolor{uefsForest}
  \textcolor{white}{\bfseries Categoria Legal} & \textcolor{white}{\bfseries Parâmetro Oficial} & \textcolor{white}{\bfseries Faixa / Percentual Mínimo (Lei nº 12.651/2012)} \\
  \midrule
  APP, Curso d'água natural (Art. 4º, I, a partir da calha regular) & Largura $< 10\ \text{m}$ \newline $10\ \text{m}$ a $50\ \text{m}$ \newline $50\ \text{m}$ a $200\ \text{m}$ \newline $200\ \text{m}$ a $600\ \text{m}$ \newline $> 600\ \text{m}$ & \textbf{$30\ \text{m}$} em cada margem \newline \textbf{$50\ \text{m}$} em cada margem \newline \textbf{$100\ \text{m}$} em cada margem \newline \textbf{$200\ \text{m}$} em cada margem \newline \textbf{$500\ \text{m}$} em cada margem \\
  \rowcolor{bgGray}
  APP, Outras feições (Art. 4º) & Nascentes / olhos d'água perenes \newline Encostas íngremes \newline Veredas (faixa marginal) & Raio mínimo de \textbf{$50\ \text{m}$} \newline Declividade superior a \textbf{$45^\circ$} ($> 100\%$) \newline Largura mínima de \textbf{$50\ \text{m}$} a partir do espaço brejoso \\
  Reserva Legal, RL (Art. 12) & Amazônia Legal, Florestas \newline Amazônia Legal, Cerrado \newline Demais regiões (incl. Caatinga/BA) & \textbf{$80\%$} da área do imóvel rural \newline \textbf{$35\%$} da área do imóvel rural \newline \textbf{$20\%$} da área do imóvel rural \\
  \rowcolor{bgGray}
  Licenciamento Trifásico (Res. CONAMA 237/1997) & \textbf{LP} (Licença Prévia) \newline \textbf{LI} (Licença de Instalação) \newline \textbf{LO} (Licença de Operação) & Aprova localização/viabilidade (fase de EIA/RIMA). \newline Autoriza início das obras conforme Plano Básico Ambiental (PBA) e controles. \newline Autoriza operação após checar cumprimento das condicionantes. \\
  \bottomrule
\end{tabularx}
\end{center}

% ============================================================================
\clearpage
\section*{2. Glossário Técnico Alfabético de Análise Ambiental}
\addcontentsline{toc}{section}{2. Glossário Técnico Alfabético de Análise Ambiental}
% ============================================================================

\begin{multicols}{2}
\small
\setlength{\parskip}{4pt}

\noindent\textbf{\color{uefsNavy}Ambiente:} Totalidade complexa, dinâmica e multiescalar resultante das interações entre componentes físico-bióticos, socioespaciais, técnicos, econômicos e institucionais.

\noindent\textbf{\color{uefsNavy}APP (Área de Preservação Permanente):} Área protegida nos termos da Lei nº 12.651/2012, coberta ou não por vegetação nativa, com a função ambiental de preservar recursos hídricos, paisagem, estabilidade geológica, biodiversidade e bem-estar humano.

\noindent\textbf{\color{uefsNavy}Aspecto Ambiental:} Elemento das atividades, produtos ou serviços de um empreendimento que pode interagir com o meio ambiente, atuando como mecanismo causador de impacto (Sánchez, 2013).

\noindent\textbf{\color{uefsNavy}Assoreamento:} Deposição excessiva de sedimentos erodidos nas vertentes no leito de rios, lagoas e reservatórios, reduzindo a calha, a capacidade de armazenamento e a qualidade da água.

\noindent\textbf{\color{uefsNavy}Avaliação de Impacto Ambiental (AIA):} Instrumento preventivo de política e gestão ambiental destinado a identificar, prever, valorar e mitigar as consequências ambientais de projetos, planos ou programas antes da tomada de decisão.

\noindent\textbf{\color{uefsNavy}Bacia Hidrográfica:} Unidade espacial delimitada por divisores topográficos que drena água, sedimentos e solutos para um único ponto de saída (exutório), constituindo a célula fundamental do planejamento ambiental integrado.

\noindent\textbf{\color{uefsNavy}Bioengenharia de Água e Solo:} Conjunto de técnicas que combinam materiais vegetais vivos (estacas, mudas, sementes) e estruturas inertes leves ou biodegradáveis para estabilizar taludes, margens e canais erosivos.

\noindent\textbf{\color{uefsNavy}Biomanta:} Manta confeccionada com fibras naturais biodegradáveis (coco, sisal, palha ou juta) fixada sobre solo exposto para dissipar a energia das gotas de chuva, reter umidade e favorecer o estabelecimento da vegetação.

\noindent\textbf{\color{uefsNavy}Coesão Radicular ($c_r$):} Acréscimo de resistência ao cisalhamento conferido ao solo pela malha de raízes finas que trabalham sob tração, elevando a estabilidade mecânica de taludes.

\noindent\textbf{\color{uefsNavy}Conectividade Hidrossedimentológica:} Grau de acoplamento físico entre fontes produtoras de escoamento e sedimentos nas vertentes e a rede de drenagem fluvial a jusante.

\noindent\textbf{\color{uefsNavy}Degradação Ambiental:} Alteração adversa das características físicas, químicas ou biológicas do ambiente que reduz sua produtividade, diversidade ou capacidade de prestar serviços ecossistêmicos.

\noindent\textbf{\color{uefsNavy}Ecodinâmica:} Abordagem sistêmica proposta por Jean Tricart (1977) que classifica as unidades da paisagem em meios estáveis, intergrades e fortemente instáveis segundo o balanço entre pedogênese e morfogênese.

\noindent\textbf{\color{uefsNavy}EIA/RIMA:} Estudo de Impacto Ambiental (documento técnico multidisciplinar completo) e seu respectivo Relatório de Impacto Ambiental (síntese acessível à sociedade), exigidos para atividades de significativa degradação (Res. CONAMA 001/1986).

\noindent\textbf{\color{uefsNavy}Erosão Laminar:} Remoção relativamente uniforme de finas camadas superficiais de solo pelo escoamento difuso em lençol; embora pouco visível no início, acarreta severa perda de fertilidade.

\noindent\textbf{\color{uefsNavy}Estaca Viva:} Estaca lenhosa de espécie com alta capacidade de brotação vegetativa cravada diretamente em taludes ou margens para enraizar, reforçar o solo e recompor a cobertura vegetal.

\noindent\textbf{\color{uefsNavy}Fragilidade Ambiental:} Grau de suscetibilidade do ambiente a processos degradacionais; divide-se em \textit{Potencial} (fatores naturais: relevo, solo, clima) e \textit{Emergente} (incorporação do uso e cobertura da terra, segundo Ross, 1994).

\noindent\textbf{\color{uefsNavy}Geossistema:} Unidade espacial dinâmica que integra potencial ecológico (abiótico), exploração biológica (biótico) e ação antrópica em interação contínua de matéria e energia (Bertrand, 2004).

\noindent\textbf{\color{uefsNavy}Geotêxtil:} Material têxtil permeável (natural ou sintético) empregado em engenharia e bioengenharia para filtração, drenagem, separação de camadas e proteção superficial contra erosão.

\noindent\textbf{\color{uefsNavy}Impacto Ambiental:} Qualquer alteração das propriedades físicas, químicas ou biológicas do meio ambiente resultante de atividades humanas, podendo ser adversa (negativa) ou benéfica (positiva).

\noindent\textbf{\color{uefsNavy}Indicador Ambiental:} Parâmetro mensurável ou índice agregado selecionado para sintetizar, monitorar e comunicar o estado, as pressões, os impactos ou as respostas de um sistema ambiental.

\noindent\textbf{\color{uefsNavy}Licenciamento Ambiental:} Procedimento administrativo pelo qual o órgão ambiental competente licencia a localização, instalação, ampliação e operação de empreendimentos potencialmente poluidores (LP, LI e LO).

\noindent\textbf{\color{uefsNavy}Morfopedologia:} Estudo integrado das relações genéticas e funcionais entre as formas do relevo, os materiais de origem, os solos e a dinâmica hídrica ao longo das vertentes.

\noindent\textbf{\color{uefsNavy}Ordenamento Territorial:} Ação normativa e estratégica do poder público voltada à organização espacial dos usos, ocupações e funções do território conforme aptidões e restrições socioambientais.

\noindent\textbf{\color{uefsNavy}Paliçada Permeável:} Barreira transversal construída com madeira, bambu ou galharia no interior de sulcos e ravinas para desacelerar o escoamento, reter sedimentos em degraus e permitir a drenagem da água sem sobrepressão.

\noindent\textbf{\color{uefsNavy}PEIR (Modelo):} Matriz conceitual Pressão--Estado--Impacto--Resposta adotada pela OCDE e pelo IBGE para estruturar indicadores ambientais de forma causal e orientada à tomada de decisão.

\noindent\textbf{\color{uefsNavy}Piping (Erosão Interna):} Remoção subsuperficial de partículas do solo pelo fluxo concentrado de água subterrânea em dutos, provocando colapso de paredes e avanço remontante de voçorocas.

\noindent\textbf{\color{uefsNavy}PRAD:} Plano de Recuperação de Áreas Degradadas; instrumento técnico-legal que estabelece diagnóstico do passivo, medidas corretivas, cronograma, orçamento e metas monitoráveis de restauração.

\noindent\textbf{\color{uefsNavy}Qualidade do Solo:} Capacidade funcional do solo de sustentar produtividade biológica, regular fluxos hídricos e manter a qualidade ambiental nos limites do ecossistema (Doran; Parkin, 1994).

\noindent\textbf{\color{uefsNavy}Ravina:} Incisão erosiva linear profunda (geralmente $> 0,5\ \text{m}$) gerada pela concentração do escoamento superficial, não mais eliminável por preparo agrícola convencional, sem ainda interceptar o lençol freático.

\noindent\textbf{\color{uefsNavy}Reserva Legal (RL):} Parcela do imóvel rural que deve ser mantida com cobertura de vegetação nativa para assegurar uso econômico sustentável, conservação da biodiversidade e abrigo de fauna (20\%, 35\% ou 80\%, conforme região/bioma).

\noindent\textbf{\color{uefsNavy}Resiliência Ambiental:} Capacidade de um sistema socioambiental absorver distúrbios, manter suas funções essenciais e reorganizar-se após perturbações naturais ou antrópicas.

\noindent\textbf{\color{uefsNavy}Risco Ambiental:} Probabilidade de ocorrência de danos humanos, materiais ou ecológicos em decorrência da interação entre uma ameaça/perigo ($P$) e a vulnerabilidade ($V$) de elementos expostos ($R = P \times V$).

\noindent\textbf{\color{uefsNavy}Topossequência:} Sequência ordenada de solos distribuídos ao longo de um perfil de vertente (do topo ao fundo de vale), refletindo variações de declividade, drenagem e espessura do manto pedológico.

\noindent\textbf{\color{uefsNavy}USLE:} Equação Universal de Perda de Solo ($A = R \cdot K \cdot L \cdot S \cdot C \cdot P$), modelo empírico utilizado para estimar a taxa média anual de erosão hídrica laminar e em sulcos.

\noindent\textbf{\color{uefsNavy}Voçoroca:} Estágio mais severo da erosão linear, caracterizado por incisões profundas que atingem o lençol freático, combinando escoamento superficial, \textit{piping} e solapamento de taludes.

\noindent\textbf{\color{uefsNavy}Vulnerabilidade:} Grau de suscetibilidade física, social, econômica ou institucional de uma população, infraestrutura ou ecossistema sofrer danos diante de uma ameaça ambiental.

\noindent\textbf{\color{uefsNavy}ZEE (Zoneamento Ecológico-Econômico):} Instrumento da Política Nacional do Meio Ambiente que compartimenta o território em zonas segundo potencialidades e fragilidades para subsidiar políticas públicas e o ordenamento territorial.
\end{multicols}

% ============================================================================
\clearpage
\section*{3. Referências Bibliográficas Oficiais da Disciplina CHF662}
\addcontentsline{toc}{section}{3. Referências Bibliográficas Oficiais da Disciplina CHF662}
% ============================================================================

\noindent Relação integral das \textbf{33 referências oficiais} constantes do Plano de Disciplina \textbf{CHF662, Análise Ambiental} (DCHF/UEFS, 2026/2), seguida das normas jurídicas complementares citadas ao longo da apostila.

\vspace{6pt}
\subsection*{Referências Oficiais do Plano de Disciplina CHF662}

\begin{list}{}{%
  \setlength{\leftmargin}{1.6em}%
  \setlength{\itemindent}{-1.6em}%
  \setlength{\itemsep}{4.5pt}%
  \setlength{\parsep}{0pt}%
  \setlength{\topsep}{4pt}%
}
\small
  \item ARRUDA, A. G. F. ``Planejamento territorial'' e ``ordenamento territorial'': uma busca da compreensão usual e epistemológica na gestão do território. \textit{Akrópolis}, Umuarama, v. 21, n. 2, p. 125--132, jul./dez. 2013.

  \item BERTALANFFY, Ludwig von. \textit{Teoria geral dos sistemas}: fundamentos, desenvolvimento e aplicações. 5. ed. Petrópolis: Vozes, 2010.

  \item BERTRAND, Georges. Paisagem e geografia física global: esboço metodológico. \textit{RA'E GA: O Espaço Geográfico em Análise}, Curitiba, n. 8, p. 141--152, 2004. DOI: \texttt{10.5380/raega.v8i0.3389}.

  \item BOTELHO, Rosangela Garrido Machado; SILVA, Antônio Soares da. Bacia hidrográfica e qualidade ambiental. In: VITTE, Antônio Carlos; GUERRA, Antônio José Teixeira (org.). \textit{Reflexões sobre a geografia física no Brasil}. Rio de Janeiro: Bertrand Brasil, 2004.

  \item BRASIL. Lei nº 6.938, de 31 de agosto de 1981. Dispõe sobre a Política Nacional do Meio Ambiente. Brasília, DF: Presidência da República, 1981.

  \item BRASIL. Lei nº 12.651, de 25 de maio de 2012. Dispõe sobre a proteção da vegetação nativa. Brasília, DF: Presidência da República, 2012.

  \item BRASIL. Instituto Brasileiro de Geografia e Estatística. \textit{Análise ambiental e gestão do território}: contribuições teórico-metodológicas. Organização de Fabio Giusti Azevedo de Britto, Letícia de Carvalho Giannella e Rogério dos Santos Seabra. Rio de Janeiro: IBGE, 2018.

  \item BRASIL. Ministério do Meio Ambiente. \textit{Procedimentos de licenciamento ambiental do Brasil}. Organização de Marco Aurélio Belmont e Pablo Ramos Andrade Villanueva. Brasília, DF: MMA, 2016.

  \item CASSETI, Valter. \textit{Geomorfologia}. Goiânia: UFG, 2007.

  \item CASTRO, Cleber Marques de; LEMOS, Clara Carvalho de. \textit{Planejamento ambiental}. Rio de Janeiro: Fundação Cecierj, 2016.

  \item CAVALCANTI, Agostinho Paula Brito. \textit{Métodos e técnicas da análise ambiental}: guia para estudos do meio ambiente. Teresina: UFPI/CCHL/DGH, 2006.

  \item CHRISTOFOLETTI, Antônio. \textit{Geomorfologia fluvial}. São Paulo: Edgard Blücher, 1981.

  \item DORAN, John W.; PARKIN, Timothy B. Defining and assessing soil quality. In: DORAN, John W. et al. \textit{Defining soil quality for a sustainable environment}. Madison: Soil Science Society of America, 1994. p. 3--21. DOI: \texttt{10.2136/sssaspecpub35.c1}.

  \item DREW, David. \textit{Processos interativos homem-meio ambiente}. Rio de Janeiro: Bertrand Brasil, 1989.

  \item DURLO, Miguel Antão; SUTILI, Fabrício Jaques. \textit{Bioengenharia}: manejo biotécnico de cursos de água. Porto Alegre: EST Edições, 2005.

  \item FIDALGO, Elaine Cristina Cardoso. \textit{Critérios para a análise de métodos e indicadores ambientais usados na etapa de diagnóstico de planejamentos ambientais}. 2003. 239 f. Tese (Doutorado em Engenharia Agrícola), Faculdade de Engenharia Agrícola, Universidade Estadual de Campinas, Campinas, 2003.

  \item GRAY, Donald H.; SOTIR, Robbin B. \textit{Biotechnical and soil bioengineering slope stabilization}: a practical guide for erosion control. New York: John Wiley \& Sons, 1996.

  \item GUERRA, Antônio José Teixeira. Geomorfologia e planejamento ambiental: conceitos e aplicações. \textit{Revista de Geografia}, Recife, v. 35, n. 4, p. 1--23, 2018.

  \item GUERRA, Antônio José Teixeira; CUNHA, Sandra Baptista da (org.). \textit{Geomorfologia e meio ambiente}. Rio de Janeiro: Bertrand Brasil, 1996.

  \item KARLEN, Douglas L. et al. Soil quality: a concept, definition, and framework for evaluation. \textit{Soil Science Society of America Journal}, Madison, v. 61, n. 1, p. 4--10, 1997. DOI: \texttt{10.2136/sssaj1997.03615995006100010001x}.

  \item LEFF, Enrique. \textit{Epistemologia ambiental}. São Paulo: Cortez, 2001.

  \item MENDONÇA, Francisco. Geografia socioambiental. \textit{Terra Livre}, São Paulo, n. 16, p. 139--158, 2001.

  \item OECD. \textit{Environmental indicators}: development, measurement and use. Paris: Organisation for Economic Co-operation and Development, 2003.

  \item PORTO-GONÇALVES, Carlos Walter. \textit{Os (des)caminhos do meio ambiente}. 14. ed. São Paulo: Contexto, 2006.

  \item RODRIGUEZ, José Manuel Mateo; SILVA, Edson Vicente da; CAVALCANTI, Agostinho Paula Brito. \textit{Geoecologia das paisagens}: uma visão geossistêmica da análise ambiental. 5. ed. Fortaleza: Edições UFC, 2017.

  \item ROSS, Jurandyr Luciano Sanches. Análise empírica da fragilidade dos ambientes naturais e antropizados. \textit{Revista do Departamento de Geografia}, São Paulo, n. 8, p. 63--74, 1994. DOI: \texttt{10.7154/RDG.1994.0008.0006}.

  \item ROSS, Jurandyr Luciano Sanches. \textit{Ecogeografia do Brasil}: subsídios para planejamento ambiental. São Paulo: Oficina de Textos, 2006.

  \item SÁNCHEZ, Luis Enrique. \textit{Avaliação de impacto ambiental}: conceitos e métodos. 2. ed. São Paulo: Oficina de Textos, 2013.

  \item SCHIECHTL, Hugo Meinhard; STERN, Roland. \textit{Water bioengineering techniques for watercourse bank and shoreline protection}. Oxford: Blackwell Science, 1996.

  \item TAUK-TORNISIELO, Sâmia Maria. \textit{Análise ambiental}: uma visão multidisciplinar. 2. ed. São Paulo: Editora UNESP, 1996.

  \item TRICART, Jean. \textit{Ecodinâmica}. Rio de Janeiro: Diretoria Técnica da SUPREN/IBGE, 1977.

  \item VEZZANI, Fabiane Machado; MIELNICZUK, João. Uma visão sobre qualidade do solo. \textit{Revista Brasileira de Ciência do Solo}, Viçosa, v. 33, n. 4, p. 743--755, 2009. DOI: \texttt{10.1590/S0100-06832009000400001}.

  \item VITAL, Marcos H. F.; CARVALHO, Frederico C.; PINTO, Marco A. C. \textit{Metodologias e técnicas para análise ambiental de projetos de investimento}. Rio de Janeiro: BNDES, 2020.
\end{list}

\vspace{8pt}
\subsection*{Normas Jurídicas Complementares Citadas na Apostila}

\begin{list}{}{%
  \setlength{\leftmargin}{1.6em}%
  \setlength{\itemindent}{-1.6em}%
  \setlength{\itemsep}{4.5pt}%
  \setlength{\parsep}{0pt}%
  \setlength{\topsep}{4pt}%
}
\small
  \item BRASIL. Lei nº 9.605, de 12 de fevereiro de 1998. Dispõe sobre as sanções penais e administrativas derivadas de condutas e atividades lesivas ao meio ambiente, e dá outras providências (Lei de Crimes Ambientais). Brasília, DF: Presidência da República, 1998.

  \item BRASIL. Conselho Nacional do Meio Ambiente (CONAMA). Resolução nº 001, de 23 de janeiro de 1986. Dispõe sobre critérios básicos e diretrizes gerais para a avaliação de impacto ambiental. Brasília, DF: \textit{Diário Oficial da União}, 17 fev. 1986.

  \item BRASIL. Conselho Nacional do Meio Ambiente (CONAMA). Resolução nº 237, de 19 de dezembro de 1997. Dispõe sobre a revisão e complementação dos procedimentos e critérios utilizados para o licenciamento ambiental. Brasília, DF: \textit{Diário Oficial da União}, 22 dez. 1997.
\end{list}
